\chapter{习题参考答案}\label{app:answers}

本书习题多为开放题：答案不止一个，评分对象是
\keyw{方法与论证}，不是结论本身。本附录给出各题的
\keyw{答题要点}与\keyw{评分点}，供自答后对照、
或供施教者批改参考。凡要求读者自选实例的题目，
答案只给判定标准；凡有计算成分的题目，给出完整过程。
建议次序：先自答，再对照；若你的答案与本附录冲突，
先检查方法，再检查结论——冲突本身可能就是
一道新的习题。

\section{第一章}\label{app:ans:01}

\begin{enumerate}[label=\textbf{1.},leftmargin=2.6em]
  \item \textbf{要点}：三维各写一句——内容（它说了什么、
        主角是谁、发生了什么事）；形式（怎么呈现：模板、
        音轨、排版、结尾套路）；立场（对谁说、以什么姿态说：
        自嘲 / 圈内认证 / 反讽评级）。“硬约定”维的判别标志：
        \keyw{用错会被社群识别为外行}（如衍生语配错场景、
        缺失结尾集体大笑）。多数梗的硬维在形式或立场，
        内容维可自由发挥——但归属取决于所选梗。
        \textbf{评分点}：能否举出“用错会被识别”的具体证据，
        而非硬维的归属结论本身。
\end{enumerate}

\begin{enumerate}[label=\textbf{2.},leftmargin=2.6em]
  \item \textbf{要点}：压力源假设指标示例——同一批账号在
        考试周内奶蛙类转发占其全部转发的比例，相对平时周的
        上升幅度（单一比例，可计数）。算法假设指标示例——
        同期话题曝光总量（平台播放量、话题页浏览）的增幅。
        若两者方向相反（曝光大增而个体转发占比不动）：
        总热度由供给侧与推荐侧驱动，使用侧的压力选择机制
        未获得支持——\keyw{竞争性解释获得相对支持}。
        \textbf{评分点}：指标是单一可计数数字；能说出
        “方向相反不是证伪，而是相对支持”这一逻辑关系。
\end{enumerate}

\begin{enumerate}[label=\textbf{3.},leftmargin=2.6em]
  \item \textbf{要点}：如全网粉丝 3800 万（《人民日报》口径）
        与 3900 万（另一来源口径）。并列而非取均值的理由：
        两数字的统计时点、口径边界（是否含海外、是否含子账号）
        可能不同；均值会制造出\keyw{任何一方都不曾主张的假口径}。
        对引用者的影响：引用时必须携带“来源 + 日期”，
        不得只引数字；也不得声称“约 3850 万”——
        并列制度把不确定性显式交还给引用者。
        \textbf{评分点}：能说出“均值 = 新造口径”这一核心。
\end{enumerate}

\section{第二章}\label{app:ans:02}

\begin{enumerate}[label=\textbf{1.},leftmargin=2.6em]
  \item \textbf{要点}：法律层面实例——官方关闭评论区、
        批量拉黑（\conf{C} 级）、法务团队与快速维权机制
        （\conf{A} 级）。意义层面实例——二创者反讽征用正典
        桥段（“糖”评级）、“潘多拉奶盒”叙事宣布脱离正典。
        两者不总是一致：法律层面的支配关系持续存在，
        意义层面的流通已不服从该边界（第九章命题
        \ref{prop:09:dual}）。
        \textbf{评分点}：两实例分属两层面；能指出
        “不一致正是迷因生态缝隙”。
\end{enumerate}

\begin{enumerate}[label=\textbf{2.},leftmargin=2.6em]
  \item \textbf{要点}：最易贴标签的通常是“文本盗猎”
        （答“抽象文化”且自圆亦可）。以文本盗猎为例，
        使用前必答的三个结构问题：(1) 盗猎\keyw{什么}——
        被挪用的文本单元（形象 / 桥段 / 音轨）；(2) \keyw{怎么}盗猎
        ——再创作的技术路径与门槛；(3) 盗猎后\keyw{功能}的变化
        ——意义、受众、流通结构各变成什么。
        \textbf{评分点}：三问缺一即视为贴标签。
\end{enumerate}

\begin{enumerate}[label=\textbf{3.},leftmargin=2.6em]
  \item \textbf{要点}：按第二章定义的三项标记逐条判定。
        (a) 讣告体作业封面：命中语体戏仿（庄重语体承载
        校园小事）+ 意义悬置（框架与对象严重错配），
        判“属于”。(b) 画质劣化但语义完整：意义未悬置；
        能否命中取决于劣化是否构成对“精致规范”的
        可辨认姿态——判“边缘”，关键看能否指认戏仿对象。
        (c) 仅画质劣化、语义清晰：三项皆不命中，
        判“不属于”。
        \textbf{评分点}：逐条对标记作答；禁止以
        “感觉很抽象”代替判定。
\end{enumerate}

\section{第三章}\label{app:ans:03}

\begin{enumerate}[label=\textbf{1.},leftmargin=2.6em]
  \item \textbf{要点}：第三季 2025-01 上线与“一个月破亿”
        有机构报道支撑（\conf{B}）；第二季播出信息仅见于
        第一版教材口径（\conf{C}）。同系列不同年份可信度
        不同，体现“\keyw{级别跟随陈述}”：等级属于
        具体陈述（某季某时点的报道是否存在），
        不属于品牌整体。
        \textbf{评分点}：能复述“级别跟随陈述而非来源”。
\end{enumerate}

\begin{enumerate}[label=\textbf{2.},leftmargin=2.6em]
  \item \textbf{要点}：层 1（用户年龄）——平台年报中
        奶龙内容 20 岁以上观众占比的时间序列；层 2
        （二创作者构成）——抽样二创作者并检索其公开简介中
        高校关键词的计数占比；层 3（场景迁移）——门店
        非亲子时段客群观察、校园立牌出现频率等自建观察。
        最难测的是层 3：场景是弥散的线下行为，
        无现成平台数据，任何代理指标都离行为更远。
        \textbf{评分点}：承认所有指标都是代理；
        层 3 指标不得依赖自陈。
\end{enumerate}

\begin{enumerate}[label=\textbf{3.},leftmargin=2.6em]
  \item \textbf{要点}：不矛盾。法务 + 维权（\conf{A}）
        是授权体系的\keyw{前置条件}（第九章）；
        关评拉黑（\conf{C}）针对“丑化二创”这一失控流通。
        两者共同服务于“形象可租赁”。依赖 C 级的部分：
        关评拉黑的时间、规模、动机细节；只依赖 A 级的部分：
        维权机制的存在及其产业逻辑。
        \textbf{评分点}：能把结论拆开，分别交代证据级别。
\end{enumerate}

\section{第四章}\label{app:ans:04}

\begin{enumerate}[label=\textbf{1.},leftmargin=2.6em]
  \item \textbf{要点}：第一句是断言形态（要求起源已证：
        载体可复核、证据链闭合）；第二句是考据形态
        （指向 + 限定：指向存在，载体不可复核须写明）。
        第一句当前不成立：关键智能体已无法访问，
        任何第三方无法复核，材料只支持“指向”，
        不支持“来自”。
        \textbf{评分点}：能说出“不可复核”在此的作用位置。
\end{enumerate}

\begin{enumerate}[label=\textbf{2.},leftmargin=2.6em]
  \item \textbf{要点}：学生自选实例。三阶段划分须正确：
        误指（被广泛指认的错误来源）、线索（推动定位的
        可观察痕迹）、发现（可复核的载体）。若该梗起源
        已公认，须指出跳过的阶段及原因（如起源系官方公布，
        则“误指”阶段无生长空间）。
        \textbf{评分点}：阶段划分与该梗实际考据史相符。
\end{enumerate}

\begin{enumerate}[label=\textbf{3.},leftmargin=2.6em]
  \item \textbf{要点}：检查程序示例。第一步：圈出修订稿中
        一切涉及具体个人的句子（关键词：本人、作者、
        当事人、其账号等）。第二步：逐句套用句式模板——
        允许：“社群中存在关于某人的说法”“该说法无法证实”；
        禁止：转述具体内容的从句（境况、去向、私生活细节）。
        第三步：核对所有“说明存在”句是否只含出处与
        存在声明。违例句数为零则通过。
        \textbf{评分点}：程序能在五分钟内执行
        （关键词清单 + 句式模板两件套齐全）。
\end{enumerate}

\section{第五章}\label{app:ans:05}

\begin{enumerate}[label=\textbf{1.},leftmargin=2.6em]
  \item \textbf{要点}：(a) 高数课本表情包：特征集保留（C1）、
        谱系同源（C2）、可识别（C3），同簇；决定性标准 C1。
        (b) 黄色大笑正常体型龙：C1 部分命中（黄、大笑），
        但若无畸变谱系出身，C2 失败，不同簇；决定性标准 C2。
        (c) 异形象同音轨：C1 失败、C3 依赖场景；若社群在
        奶蛙语境中稳定使用，可划入“形式衍生”边缘。
        \textbf{评分点}：本题意图在让读者发现三条标准
        可能给出不一致判定——\keyw{承认模糊比强行裁决
        更符合第五章立场}，能写出这一层为优。
\end{enumerate}

\begin{enumerate}[label=\textbf{2.},leftmargin=2.6em]
  \item \textbf{要点}：按第五章记号，感染变化率
        $\mathrm{d}I/\mathrm{d}t = I \times
        (\text{接触项} - \text{恢复项})$。$\omega = 0$ 时
        $S(t)$ 单调不增，故接触项随时间不增；恢复项不增减
        （或单调）。令 $\mathrm{d}I/\mathrm{d}t = 0$：
        等式至多在一点成立（或成立为一个区间，对应平台期）。
        因此 $I(t)$ 先增、后（至多）减，至多一个极大值。
        若 $S(t)$ 始终高于阈值，则 $I$ 单调增——峰在
        观测窗之外，这正是命题中“\keyw{至多}”一词的含义。
        \textbf{评分点}：用到“$S$ 单调”与“驻点条件”
        两个要素；“至多”的边界情形。
\end{enumerate}

\begin{enumerate}[label=\textbf{3.},leftmargin=2.6em]
  \item \textbf{要点}：判据示例——先作 7 日平滑；
        峰 = 局部极大且两侧回落超过峰值的三成；
        双峰 = 两峰间隔不小于两周且谷值低于两峰均值六成。
        稳健化：对平滑窗宽与谷值比例做扰动，形态不变者
        才报告。不能识别的：无法区分“同一梗的二次爆发”
        与“两个独立梗的叠加”；无法识别免疫衰减参数本身；
        无法排除平台限流造成的假谷。
        \textbf{评分点}：判据可执行 + 至少一条
        “本方案不能识别什么”。
\end{enumerate}

\section{第六章}\label{app:ans:06}

\begin{enumerate}[label=\textbf{1.},leftmargin=2.6em]
  \item \textbf{要点}：逐条对照 (C1)(C2)(C3)，指出
        起决定作用的一条。主张移出某变种者，必须指认
        其丢失了哪个“不可再丢”的特征（如只剩黄色、
        失去大笑与谱系）。无标准答案。
        \textbf{评分点}：判定过程可复核（三条标准逐一给出
        满足 / 不满足及证据）。
\end{enumerate}

\begin{enumerate}[label=\textbf{2.},leftmargin=2.6em]
  \item \textbf{要点}：以“典”为例——判定操作：
        复读已成套路的论证模板（可观察特征：句式雷同、
        引用链封闭）。被官方或广告征用后的变化：
        评级权从社群转移到被评对象，词义“通货膨胀”
        （什么都典 = 无物可典），社群或弃用、或转义为
        更深的反讽。
        \textbf{评分点}：把“效力变化”操作化
        （使用场景的收敛或迁移），不停留在“变味了”。
\end{enumerate}

\begin{enumerate}[label=\textbf{3.},leftmargin=2.6em]
  \item \textbf{要点}：行为指标示例——统计转发该文本的用户
        在其后一段时间内是否仍参与官方相关负面二创；
        文本指标示例——该文本变体的演化方向（强化立场、
        还是滑向纯玩梗）。不一致时\keyw{接受行为指标}：
        文本只能证明“社群愿意如此自我叙述”（自我理解），
        不能证明个体按此行动（动机证据）——这正是第六章
        对“奶蛙不叛”的处理方式。
        \textbf{评分点}：说出两指标的分界理由，
        且裁决与第六章定位一致。
\end{enumerate}

\section{第七章}\label{app:ans:07}

\begin{enumerate}[label=\textbf{1.},leftmargin=2.6em]
  \item \textbf{要点}：四列齐全（压力源—补偿行为—
        可观察指标—竞争性解释）；指标可在不接触隐私的
        前提下计数（公开帖计数、话题标签频次）；
        证伪条件具体到观察层面。
        \textbf{评分点}：自查一步（指标是否隐私安全）
        必须真的执行——若原始指标涉隐私且未改写，本题不及格。
\end{enumerate}

\begin{enumerate}[label=\textbf{2.},leftmargin=2.6em]
  \item \textbf{要点}：判别实验示例——把“圈内身份”与
        “压力水平”正交：高压力圈外组（临近大考的圈外学生）
        与低压力圈内组（假期中的同好会核心成员）。转发频率
        随压力变 → 支持群体稀释；随身份变 → 支持
        身份认证。自然实验版本：考试周前后同一人群的
        频率变化（压力变、身份不变）。
        \textbf{评分点}：两组至少在一个关键变量上正交；
        说清什么结果对应哪个解释。
\end{enumerate}

\begin{enumerate}[label=\textbf{3.},leftmargin=2.6em]
  \item \textbf{要点}：若 M1 证伪而 M3 独立成立：现象
        可持续性\keyw{上升}——情绪压力是周期性驱动
        （考试周），风格规范是长期性的（社群身份），
        现象从“周期性情绪出口”转为“稳定风格规范”，
        寿命预期延长、季节波动减弱。反之（仅 M1 成立）：
        现象随压力节律波动、随毕业潮流失受体，
        寿命与校园节律绑定。
        \textbf{评分点}：两个方向都写；把“分层可独立证伪”
        与“寿命预测”连接起来。
\end{enumerate}

\begin{enumerate}[label=\textbf{4.},leftmargin=2.6em]
  \item \textbf{要点}：策略选项示例。(1) 转移视线：
        正面内容供给挤压丑化空间——改变 M4 的场景供给结构；
        (2) 官方下场玩梗：把 M5 的逆反能量转化为授权流量，
        风险是触发 M3 的反精致反弹（收编指控）；
        (3) 选择性维权：对商用侵权维权、对玩梗沉默——
        把 M5 的“认证信号”与“法律信号”分离；
        (4) 开放共创授权：重设 M5 的权利结构，
        使打压—认证回路失去输入。
        \textbf{评分点}：每个选项对应到具体机制的变动，
        并给出副作用。
\end{enumerate}

\section{第八章}\label{app:ans:08}

\begin{enumerate}[label=\textbf{1.},leftmargin=2.6em]
  \item \textbf{要点}：逐位对齐（两词均 5 音节）：
        嘎 / 刚——声母 \texttt{g} 同（0）、韵母
        \texttt{a}/\texttt{ang}（1）；嘎 / 刚（1）；
        滴 / 点——声母 \texttt{d} 同（0）、韵母
        \texttt{i}/\texttt{ian}（1）；辣 / 了——声母
        \texttt{l} 同（0）、韵母 \texttt{a}/\texttt{e}（1）；
        虾 / 谁——声母 \texttt{x}/\texttt{sh}（1）+
        韵母 \texttt{ia}/\texttt{ui}（1）= 2。
        总代价 6 分，$d = 6/5 = 1.20$。逐位对齐即最优：
        任何增删（每次 1.5）与重排的代价都不低于它。
        仍可成立的理由：等价关系由三条标准\keyw{联合}判定
        ——(2) 场景稳定（该空耳在特定桥段被社群稳定使用）、
        (3) 实例可检索（存在公开二创实例）；$d$ 只服务
        标准 (1) 的\keyw{排序}，而标准 (1) 的文字表述是
        “大体对应”，无数值阈值——音节数一致、叠音结构
        对应（嘎嘎$\leftrightarrow$刚刚）已满足其质性要求。
        \textbf{评分点}：对齐方案 + 计算正确；
        能引用 (2)(3) 与“$d$ 是排序指标”两条警告。
\end{enumerate}

\begin{enumerate}[label=\textbf{2.},leftmargin=2.6em]
  \item \textbf{要点}：最小对照示例。样本：同一平台、
        同一时间窗内“改结尾”与“保留结尾”的奶蛙二创视频
        各 $N$ 条（题材配对）。变量：是否改结尾（二值）、
        传播表现（转发 / 点赞，取对数）。统计量：两组中位数
        比较或配对差值检验。混杂排除：只取作者
        \keyw{首次发布}的作品（人人都不熟练），
        或把作者历史发布量作为配对 / 协变量。
        \textbf{评分点}：样本—变量—统计量三件套齐全；
        混杂处理具体到操作。
\end{enumerate}

\begin{enumerate}[label=\textbf{3.},leftmargin=2.6em]
  \item \textbf{要点}：机制设计示例——对单账号单位时间的
        批量低原创发布设“梯度曝光衰减”（类累进排放税：
        发布量越大，单条初始曝光越低），费率公开。
        反弹分析：社群会把限流本身再梗化
        （“被吞了”成为新素材），触发 M5 逆反认证——
        \keyw{治理动作成为新的认证信号}；同时限流损伤 M2
        的群体稀释通道，压力表达将寻找替代出口。
        \textbf{评分点}：有“费率”结构；反弹对应到
        具体机制并说明传播后果。
\end{enumerate}

\section{第九章}\label{app:ans:09}

\begin{enumerate}[label=\textbf{1.},leftmargin=2.6em]
  \item \textbf{要点}：排序示例：结尾集体大笑 $>$ 奶蛋 $>$
        品牌联名奶茶杯。依据三维：内容可换——大笑不依赖
        具体桥段，奶蛋仍绑定“奶伴”对战叙事，联名杯绑定
        品牌与产品；载体可换——大笑跨平台流通，奶蛋限于
        B 站语境，联名杯限于线下商品；署名可弃——三者
        均已匿名化，此项不构成区分。
        \textbf{评分点}：三维分别论证，不笼统排序。
\end{enumerate}

\begin{enumerate}[label=\textbf{2.},leftmargin=2.6em]
  \item \textbf{要点}：判据示例。官方视角：\keyw{可执行性}
        ——对象超出可执法清单（匿名、跨平台、线下化）
        即失控；社群视角：\keyw{伤害}——出现实质损害
        （诈骗、人身攻击）才失控。差异根源：官方以权利
        结构为界，社群以内容性质为界。
        \textbf{评分点}：判据可操作（能裁决具体案例）；
        能说明两视角的界线性质不同。
\end{enumerate}

\begin{enumerate}[label=\textbf{3.},leftmargin=2.6em]
  \item \textbf{要点}：观点一隐含前提：匿名 = 放弃权利主张
        （把“不可识别”当作“权利放弃”）。观点二隐含前提：
        社群劳动创造的价值应归属劳动者（劳动价值预设 +
        可识别的贡献分配）。E3 支持的\keyw{分析姿态}：
        叙述二创形象时附署名权归属，使创作者作为权利主体
        在文本中可见；E3 不支持的\keyw{结论}：两种观点都是
        规范性结论，E3 作为考据伦理规则不裁决分配正义
        ——它只防止权利主体在叙述中消失。
        \textbf{评分点}：分清“姿态”与“结论”两个层次。
\end{enumerate}

\section{第十章}\label{app:ans:10}

\begin{enumerate}[label=\textbf{1.},leftmargin=2.6em]
  \item \textbf{要点}：低—中强度分支：例行移除（校方 0）
        + 纪念文本（学生 1）+ 线上存续（持久物 1）。
        高强度分支：公开通报或追责（校方 2）$\to$
        组织化表态（学生 2）$\to$ 制度化对抗或申请合规化
        （持久物 2）。转折点候选：校方是否\keyw{公开追责}
        （把例行处置升格为立场表态）、学生是否出现
        \keyw{组织化}动作（从个体玩梗变为集体对抗）。
        \textbf{评分点}：两个分支的三维读数齐全；
        转折点指认到具体动作。
\end{enumerate}

\begin{enumerate}[label=\textbf{2.},leftmargin=2.6em]
  \item \textbf{要点}：观察点——流动投放点与固定立牌
        各一（可比）；时段——投放后固定窗口
        （如午间 60 分钟 $\times$ 5 天）；计数规则——
        合影率 = 驻足者中拍照人数 / 驻足总人数，
        “驻足”定义为停留不短于 5 秒；混杂——天气、
        考试周、人流构成（逐场记录作控制变量）；
        伦理——公共场所匿名计数，不摄录人脸，
        不跟踪个体。
        \textbf{评分点}：计数规则事先写定；
        至少两类混杂及处理。
\end{enumerate}

\begin{enumerate}[label=\textbf{3.},leftmargin=2.6em]
  \item \textbf{要点}：找到的文本须有可撤回装置
        （戏谑语体、夸张到不可能当真、匿名性）；
        若不成立，须指出其替代机制（实名表态、
        组织化申诉、直接协商）。加分项：把该文本与
        赞诗作结构对比（同一功能、不同装置）。
        \textbf{评分点}：“可撤回装置”被具体指认，
        而非笼统说“也是玩梗”。
\end{enumerate}

\section{第十一章}\label{app:ans:11}

\begin{enumerate}[label=\textbf{1.},leftmargin=2.6em]
  \item \textbf{要点}：前瞻替代——固定样本日记法：
        招募小样本（如 60 人），考试周前后连续 21 天
        每日推送一题（“今日是否转发奶蛙类内容 0/1 +
        情境勾选”）。取舍：日记法单样本成本高、
        流失率高，但消除近因效应与情绪一致性回忆；
        事件—时间对齐精度远高于回填。
        \textbf{评分点}：说出回填法的两类系统偏差
        及其方向；成本与效度的取舍双向呈现。
\end{enumerate}

\begin{enumerate}[label=\textbf{2.},leftmargin=2.6em]
  \item \textbf{要点}：反方最强理由：辨识“故意劣化”
        依赖整体观感（格式塔），拆成标记会系统性漏掉
        明显故意做差但无标记的内容——\keyw{效度}受损。
        正方最强理由：无标记的整体观感编码无法区分
        故意与真的做差，编码者一致性必然崩溃——
        \keyw{信度}崩塌。裁决示例：信度优先
        （宁可漏、不可滥），折中方案——标记编码为主、
        入假设检验；另设“整体判断”辅助列，不入检验，
        只用于生成后续假设。
        \textbf{评分点}：说清信度—效度冲突的结构；
        裁决有明确取舍而非和稀泥。
\end{enumerate}

\begin{enumerate}[label=\textbf{3.},leftmargin=2.6em]
  \item \textbf{要点}：修改示例——措辞去“屏蔽”改中性
        （“如果你刷奶蛙的频率变低了，你的使用会怎么变？”）；
        位置后移至访谈末尾；附正常化选项（“也许不会变
        ——很多人过段时间就换新梗了”）以降诱导。
        若修改后失去分析价值的理由：M5 检验需要
        “压制—反应”的强刺激设定，中性化的“频率变低”
        测到的只是兴趣衰退（M4 范围），无法生成
        压制反弹假设——此时应承认该题不能温和化，
        改以第十章已发生的移除事件作自然实验替代。
        \textbf{评分点}：修改真正降低诱导（三处之一即可）；
        “失去价值”的论证触及 M5 与 M4 的区分。
\end{enumerate}

\section{第十二章}\label{app:ans:12}

\begin{enumerate}[label=\textbf{1.},leftmargin=2.6em]
  \item \textbf{要点}：条目号正确；原句真实存在
        （注明章节位置）；修改句符合该条目的操作要点；
        依据引用条目原文。批判框所列三处病灶之外的自选
        例子同样有效——\keyw{发现本书未登记的新违例，
        应评为高质量答案}。
        \textbf{评分点}：四要素齐全；修改不引入新违例
        （如把 C 级句改为 A 级口吻）。
\end{enumerate}

\begin{enumerate}[label=\textbf{2.},leftmargin=2.6em]
  \item \textbf{要点}：前三步示例。(1) 更新事实手册
        （伦理检查点：重新访问 C 级来源时仍守 E1--E3）；
        (2) 获取校本材料授权（伦理检查点：当事人同意 +
        匿名化选项）；(3) 按第十章 C 级纪律写案例
        （伦理检查点：单一来源不作论证前提）。
        最易省略的是第 (2) 步——“发在群里好像就是
        公开的”。后果：校本材料成为无授权的个人信息，
        违反最小化收集；续编版本同时破坏 CC BY-NC-SA
        的署名链条。
        \textbf{评分点}：每步有检查点；对“最易省略”
        的论证触及“公开可见 $\ne$ 可自由使用”。
\end{enumerate}

\begin{enumerate}[label=\textbf{3.},leftmargin=2.6em]
  \item \textbf{要点}：反例候选——鲁迅杂文（幽默进入
        论证且增强说服力）、讽喻传统（《动物农场》以
        幽默承担政治论证）、以讽刺组织论据的学术写作。
        成立前提：(1) 幽默与论证指向同一对象；
        (2) 受众能解码幽默；(3) \keyw{幽默组织证据，
        而不替代证据}。奶蛙研究是否具备：部分具备——
        社群文本（赞诗）本身就是幽默论证，可作研究对象；
        但研究者的幽默（教材体例）与论证对象（社群行为）
        分属不同层次，若研究者的幽默直接充当论证
        （以段子代替机制分析），前提 (3) 即被破坏——
        这正是第十二章把幽默限制在体例层的理由。
        \textbf{评分点}：说出前提 (3)；能把反例前提
        与本书的分层处理对接。
\end{enumerate}
